
\uline{régulier}, \uline{si}, \uline{localement sur} @@M0@@, \uline{il existe un} @@M1@@-\uline{Module libre de type fini} @@M2@@, \uline{et un homomorphisme régulier surjectif} : @@M3@@.

\uline{Soit} @@M4@@ \uline{une immersion de préschémas}, \uline{et soit} @@M5@@ \uline{un ouvert de} @@M6@@ \uline{tel que} @@M7@@ \uline{et que} @@M8@@ \uline{soit une immersion fermée de} @@M9@@ \uline{dans} @@M10@@. \uline{On dit que} @@M11@@ \uline{est régulière si l'idéal de} @@M12@@ \uline{définissant} @@M13@@ \uline{est régulier} (\uline{condition qui ne dépend pas de} @@M14@@). \uline{Si} @@M15@@, \uline{on dit que} @@M16@@ \uline{est régulière en} @@M17@@ \uline{s'il existe un voisinage} @@M18@@ \uline{de} @@M19@@ \uline{tel que l'immersion} @@M20@@ \uline{soit régulière}. \uline{Pour qu'une immersion soit régulière}, \uline{il est donc nécessaire et suffisant qu'elle le soit au voisinage de chaque point}.

\emph{Cette définition est en conflit avec} EGA IV 16.9.2, \emph{et remplace cette dernière}. Toutefois, dans le cas noethérien, les deux définitions sont équivalentes (EGA IV 19.5.1). Une immersion qui peut localement être définie par une suite régulière est régulière; une immersion régulière est quasi-régulière (1.3).

\textbf{1.4.1.} Soient @@M21@@ un schéma, @@M22@@ un Idéal de type fini de @@M23@@, @@M24@@ un point de @@M25@@. Nous dirons qu'une famille de sections @@M26@@ de @@M27@@ sur un voisinage de @@M28@@ forme un \uline{système minimum de générateurs de} @@M29@@ \uline{au voisinage de} @@M30@@ si les @@M31@@ sont une famille de générateurs de @@M32@@ de nombre minimum, i.e. leurs images dans @@M33@@ forment une base de cet espace vectoriel sur @@M34@@. On notera que les @@M35@@ engendrent @@M36@@ sur un voisinage de @@M37@@.

\textbf{1.4.2.} Dire qu'une immersion est régulière, c'est encore dire qu'elle est définie dans un certain ouvert de @@M38@@ par un idéal de type fini @@M39@@, et que pour tout point @@M40@@, et tout système minimum de générateurs de @@M41@@ au voisinage de @@M42@@, le complexe de Koszul associé à cette suite est acyclique au voisinage de @@M43@@: soit en effet @@M44@@ un système minimum de
